% Cech descent for positive Hall halves on CY3 atlases.

\documentclass[11pt]{article}
\usepackage[localtheorems]{raeez-math-template}

\providecommand{\CC}{\mathbb{C}}
\providecommand{\ZZ}{\mathbb{Z}}
\providecommand{\cA}{\mathcal{A}}
\providecommand{\cC}{\mathcal{C}}
\providecommand{\cH}{\mathcal{H}}
\providecommand{\cM}{\mathcal{M}}
\providecommand{\cO}{\mathcal{O}}
\providecommand{\cU}{\mathcal{U}}
\providecommand{\cY}{\mathcal{Y}}
\providecommand{\Ran}{\operatorname{Ran}}
\providecommand{\Fact}{\operatorname{Fact}}
\providecommand{\Crit}{\operatorname{Crit}}
\providecommand{\CoHA}{\operatorname{CoHA}}
\providecommand{\Perf}{\operatorname{Perf}}
\providecommand{\Rep}{\operatorname{Rep}}
\providecommand{\Spec}{\operatorname{Spec}}
\providecommand{\Tot}{\operatorname{Tot}}
\providecommand{\Alg}{\operatorname{Alg}}

\newtheorem{theorem}{Theorem}[section]
\newtheorem{proposition}[theorem]{Proposition}
\newtheorem{lemma}[theorem]{Lemma}
\newtheorem{conjecture}[theorem]{Conjecture}
\theoremstyle{definition}
\newtheorem{definition}[theorem]{Definition}
\newtheorem{construction}[theorem]{Construction}
\theoremstyle{remark}
\newtheorem{remark}[theorem]{Remark}

\title{\v{C}ech descent for positive Hall halves on CY$_3$ atlases}
\author{}
\date{April 2026}

\begin{document}
\maketitle

\section{Cyclic Hall atlases}

A \v{C}ech descent theorem for compact CY$_3$ Hall algebras starts with
an atlas rather than a global quiver. For an arbitrary compact
Calabi-Yau threefold, a toric quiver presentation is extra structure:
it turns local critical CoHA calculations into a descent problem.

\begin{definition}[cyclic Hall atlas]\label{def:cyclic-hall-atlas}
Let $X$ be a smooth compact CY$_3$ with holomorphic volume form
$\Omega_X$. A cyclic Hall atlas on $X$ consists of:
\begin{enumerate}
\item a finite affine cover in the \'etale topology
$\cU=\{U_\alpha\}_{\alpha\in I}$ and all nonempty intersections
$U_{\alpha_0\cdots\alpha_p}$;
\item for each $U_{\alpha_0\cdots\alpha_p}$, an oriented $3$-Calabi-Yau
quiver-with-potential model $(Q_{\alpha_0\cdots\alpha_p},
W_{\alpha_0\cdots\alpha_p})$ presenting the chosen NCCR or tilting chart
for $\Perf(U_{\alpha_0\cdots\alpha_p})$;
\item restriction Morita equivalences between these models for every
inclusion in the \v{C}ech poset;
\item coherent comparison isomorphisms of the orientation lines and
vanishing-cycle complexes, satisfying the pentagon identity on quadruple
overlaps.
\end{enumerate}
The final clause is part of the datum. It is the derived replacement for
a single global quiver potential.
\end{definition}

For a chart $U_J=U_{\alpha_0\cdots\alpha_p}$ write
$(Q_J,W_J)$ for its quiver with potential. Its positive Hall half is
the critical CoHA
\[
\cY^+_J
\;=\;
\CoHA(Q_J,W_J)
\;=\;
\bigoplus_{\mathbf d\in\ZZ_{\geq0}^{Q_{J,0}}}
H^{\mathrm{BM}}_{*,G_{J,\mathbf d}}\!\left(
\Crit(\operatorname{Tr}W_{J,\mathbf d}),
\varphi_{\operatorname{Tr}W_{J,\mathbf d}}\otimes o_{J,\mathbf d}
\right),
\]
where $G_{J,\mathbf d}=\prod_{i\in Q_{J,0}}\mathrm{GL}_{d_i}$ and
$o_{J,\mathbf d}$ is the orientation-line choice. Multiplication is Hall
convolution along the stack of short exact sequences of representations,
with Thom-Sebastiani identifying the vanishing-cycle factors. For
$U_J\simeq\CC^3$ with the Jordan triple-loop potential
$W=\operatorname{tr}(XYZ-XZY)$, the Schiffmann-Vasserot 2013 theorem
identifies
\[
\CoHA(\CC^3)\;\simeq\;Y^+(\widehat{\mathfrak{gl}}_1)
\]
as an associative positive half. For smooth toric CY$_3$ targets whose
compact four-cycle sector is empty, Rapcak-Soibelman-Yang-Zhao 2020
identify the critical CoHA with the positive half
$Y^+(\widehat{\mathfrak g}_{Q})$ of the affine super Yangian attached to
the toric quiver.

\section{The \v{C}ech totalisation}

The atlas of Definition~\ref{def:cyclic-hall-atlas} gives a cosimplicial
diagram in the infinity category of $E_1$ algebras in complexes,
\[
[p]\longmapsto
\prod_{\alpha_0,\ldots,\alpha_p}
\cY^+_{\alpha_0\cdots\alpha_p},
\]
where the coface maps are restriction maps followed by the Morita and
vanishing-cycle comparison isomorphisms. Degeneracy maps come from
identity refinements of the intersections.

\begin{definition}[descended positive half]\label{def:descended-positive-half}
The descended positive Hall half attached to the cyclic Hall atlas
$\cU$ is the totalisation
\[
Y^+_{\cU}(X)
\;:=\;
\Tot\!\left(
\prod_{\alpha}\cY^+_{\alpha}
\rightrightarrows
\prod_{\alpha,\beta}\cY^+_{\alpha\beta}
\longrightarrow
\prod_{\alpha,\beta,\gamma}\cY^+_{\alpha\beta\gamma}
\;\cdots
\right)
\]
in $\Alg_{E_1}(\mathrm{Ch}_{\CC})$. Equivalently, it is
$R\Gamma(X,\cY^+_{\cU})$ for the sheaf convention. In the cosheaf
convention for factorisation algebras on $\Ran X$, the corresponding
factorisation object is obtained before applying sections; its derived
sections are the same totalisation.
\end{definition}

\begin{proposition}[descent of the ordered product]\label{prop:ordered-product-descent}
The object $Y^+_{\cU}(X)$ carries a canonical $E_1$ algebra structure.
If every chart CoHA carries a coproduct natural under the atlas Morita
equivalences, then $Y^+_{\cU}(X)$ is an $E_1$ bialgebra. The positive
half acquires an antipode only after Hall-Drinfeld double data.
\end{proposition}

\begin{proof}
For operadic algebras in complexes over $\CC$, the forgetful functor
from $E_1$ algebras to complexes creates limits. The cosimplicial
diagram above is a diagram of $E_1$ algebras because Hall convolution is
natural under the restriction correspondences and under the chosen
Morita equivalences. Its totalisation therefore inherits the ordered
product.

For the coproduct, the same argument applies to the diagram in
$E_1$ bialgebras once the local coproducts are natural under the
transition equivalences. The antipode belongs to a Hopf or quasi-Hopf
completion after a nondegenerate pairing and the Hall-Drinfeld double
have been supplied; it is absent from the positive half constructed by
\v{C}ech descent.
\end{proof}

\begin{remark}[full nerve versus two-term equaliser]\label{rem:full-nerve}
The two-term equaliser
\[
\operatorname{eq}\!\left(
\prod_{\alpha}\cY^+_{\alpha}
\rightrightarrows
\prod_{\alpha,\beta}\cY^+_{\alpha\beta}
\right)
\]
is valid under a split one-coskeletal cover or after higher
\v{C}ech cohomology has vanished for the coefficient diagram. The
general compact CY$_3$ atlas uses the full totalisation in
Definition~\ref{def:descended-positive-half}.
\end{remark}

\begin{proposition}[scalar characters and Hall descent]\label{prop:character-vs-algebra-descent}
Let $R$ be a commutative coefficient algebra. A compatible family of
graded characters
\[
\chi_J\colon \cY^+_J\longrightarrow R
\]
on the \v{C}ech nerve descends to a scalar character
$\chi_{\cU}\colon Y^+_{\cU}(X)\to R$. This scalar descent constructs
only the character. Algebra descent requires transition maps preserving
multiplication; bialgebra descent additionally requires coproduct
descent; Drinfeld-double descent requires a nondegenerate pairing, the
radical quotient, and the chosen completion.
\end{proposition}

\begin{proof}
Compatibility of the $\chi_J$ says exactly that the $\chi_J$ form a
map from the cosimplicial Hall diagram to the constant diagram $R$.
Totalisation gives the induced scalar map. The data discarded by
$\chi_J$ are precisely the tensors
\[
m_J\colon \cY^+_J\otimes\cY^+_J\to\cY^+_J,\qquad
\Delta_J\colon \cY^+_J\to\cY^+_J\widehat{\otimes}\cY^+_J,
\]
and, for a double, the pairing and completion. A graded character
records the graded dimension or trace of each charge sector; it does
not record the Hall correspondence defining $m_J$, the compatibility
square for $\Delta_J$, or the kernel of the pairing. Hence a glued
character is a character of the descended object only after the
corresponding algebraic descent datum has been supplied.
\end{proof}

\section{Overlap coherences}

\begin{construction}[transition data]\label{cstr:transition-data}
For a double overlap $U_{\alpha\beta}$, the atlas supplies a Morita
equivalence of oriented $3$-CY dg models
\[
M_{\alpha\beta}\colon
(\cA_\alpha,W_\alpha)|_{U_{\alpha\beta}}
\longrightarrow
(\cA_\beta,W_\beta)|_{U_{\alpha\beta}}.
\]
It induces an equivalence of critical loci and vanishing-cycle complexes,
\[
g_{\alpha\beta}\colon
\bigl(\Crit(W_\alpha),\varphi_{W_\alpha}\bigr)|_{U_{\alpha\beta}}
\longrightarrow
\bigl(\Crit(W_\beta),\varphi_{W_\beta}\bigr)|_{U_{\alpha\beta}},
\]
hence an $E_1$ algebra equivalence
$g_{\alpha\beta}\colon\cY^+_{\alpha\beta}\to\cY^+_{\beta\alpha}$.
\end{construction}

\begin{lemma}[triple-overlap coherence]\label{lem:triple-coherence}
On a triple overlap $U_{\alpha\beta\gamma}$, the descent datum is a
specified $2$-cell
\[
\eta_{\alpha\beta\gamma}\colon
g_{\beta\gamma}\circ g_{\alpha\beta}
\Longrightarrow
g_{\alpha\gamma}
\]
in the infinity groupoid of $E_1$ algebra equivalences of
$\cY^+_{\alpha\beta\gamma}$. The $2$-cells satisfy the pentagon
identity on quadruple overlaps.
\end{lemma}

\begin{proof}
Derived Morita equivalences of the chart dg algebras compose only up to
coherent natural isomorphism. Passing to critical loci and
vanishing-cycle sheaves transports these natural isomorphisms to
$2$-cells of Hall algebra equivalences. The orientation-line
compatibility in Definition~\ref{def:cyclic-hall-atlas} makes the
Thom-Sebastiani convolution products compatible with those $2$-cells.
The pentagon identity is exactly the quadruple-overlap condition in the
definition of a cyclic Hall atlas.
\end{proof}

\begin{remark}[strict global potentials]\label{rem:strict-global-potential}
The obstruction to strictifying the atlas to one global quiver
potential $(Q_X,W_X)$ is a nonabelian \v{C}ech class in the $2$-group of
oriented Morita autoequivalences. Its first-order shadow lies in the
\v{C}ech cohomology of local Hochschild cochains in degree $2$. A
nonzero obstruction blocks strict global presentation and still permits
descent in the infinity category.
\end{remark}

\section{Positive-half boundary}

\begin{theorem}[descent stops at the positive half]\label{thm:positive-half-boundary}
The object $Y^+_{\cU}(X)$ constructed from a cyclic Hall atlas is an
$E_1$ Hall positive half. It contains the chartwise positive modes and
the coherently glued Hall product. Negative halves, Cartan doubles,
universal $R$-matrices, vertex-algebra OPEs, and Borcherds-Kac-Moody
denominator identities enter through separate double, centre,
evaluation, and automorphic-product inputs.
\end{theorem}

\begin{proof}
Each entry of the \v{C}ech diagram is a critical CoHA positive half.
Limits in $\Alg_{E_1}$ preserve the operations already present in the
diagram. Negative modes and Cartan modes
enter through the Hall-Drinfeld double
\[
D\bigl(Y^+\bigr)
\;=\;
Y^-\widehat{\bowtie}Y^0\widehat{\bowtie}Y^+,
\]
which requires a nondegenerate Hopf pairing, a PBW triangular
decomposition, a radical quotient, and a completion. Braiding on
representations enters
through the Drinfeld centre
\[
\mathcal Z\!\left(\Rep^{E_1}(Y^+_{\cU}(X))\right),
\]
where half-braidings carry the $R$-matrix. A
$\mathcal W_{1+\infty}$ vertex algebra appears only after a full
Yangian or doubled algebra acts on a chosen vacuum module through an
evaluation morphism. A Borcherds-Kac-Moody algebra requires a separate
Borcherds product or denominator input, with weight
$\kappa_{\mathrm{BKM}}=c_N(0)/2$ in the Gritsenko-Borcherds
normalisation. On the primitive CHL rows $N\in\{1,2,3,4,6\}$ one has
\[
c_N(0)=(10,8,6,4,2),\qquad
\kappa_{\mathrm{BKM}}(\Phi_N)=(5,4,3,2,1).
\]
The orders $N\in\{5,7,8\}$ belong to separate Nikulin or
Gritsenko-Clery rows and carry their own Borcherds inputs.
\end{proof}

\begin{center}
\begin{tabular}{c|c|c}
Layer & Operation & Output \\
\hline
\v{C}ech Hall descent & totalisation in $\Alg_{E_1}$ &
$Y^+_{\cU}(X)$ \\
Hall double & pairing, PBW, completion &
$D(Y^+)$ \\
Drinfeld centre & half-braidings on $\Rep^{E_1}$ &
$E_2$ braided representation category \\
Vacuum evaluation & representation of the full double &
$\mathcal W$-algebra module action \\
Borcherds lift & automorphic product input &
BKM algebra and $\kappa_{\mathrm{BKM}}$
\end{tabular}
\end{center}

\section{Dimension three and scalar invariants}

The CY-to-chiral construction at dimension $d=3$ has native
$E_1$ output. The braided $E_2$ structure lives on the Drinfeld centre
of the $E_1$ representation category on the constructed loci. The
ambient $E_3$ structure of six-dimensional holomorphic Chern-Simons
local observables is a separate factorisation-algebra object; comparing
it with Hall descent requires the framed hCS-Hall datum, orientation
compatibility, compact-support realisation, and completion.

For a compact CY$_3$, Serre duality gives the Hodge-supertrace scalar
\[
\kappa_{\mathrm{ch}}^{\mathrm{Hodge}}(X)
\;=\;
\sum_{q=0}^{3}(-1)^q h^{0,q}(X)
\;=\;(h^{0,0}-h^{0,3})-(h^{0,1}-h^{0,2})
\;=\;0.
\]
Thus $\kappa_{\mathrm{cat}}(X)=\chi(\cO_X)=0$ on the categorical Euler
lane. This scalar is independent of the descended positive Hall half.
The standard local normalisations are
\[
\kappa_{\mathrm{ch}}^{\mathrm{Heis}}(\CC^3)=1,\qquad
\kappa_{\mathrm{ch}}\!\left(\Tot(K_{\mathbb{P}^2})\right)=\frac32,\qquad
\kappa_{\mathrm{ch}}(X_{\mathrm{con}})=1,
\]
where
$X_{\mathrm{con}}=\Tot(\mathcal O_{\mathbb{P}^1}(-1)\oplus
\mathcal O_{\mathbb{P}^1}(-1))$ is computed by the direct
McKay/conifold quiver calculation. It is a curve-core conifold geometry,
not a local surface. BKM weights such as
$\kappa_{\mathrm{BKM}}(\Phi_N)=c_N(0)/2$ live in the automorphic
denominator construction. None of these scalars is produced by
\v{C}ech totalisation of $Y^+$ without an explicit character map.

\begin{theorem}[one-chart check]\label{thm:one-chart-check}
For the one-chart atlas $X=\CC^3$ with the Jordan triple-loop potential,
Definition~\ref{def:descended-positive-half} gives
\[
Y^+_{\{\CC^3\}}(\CC^3)
\;\simeq\;
\CoHA(\CC^3)
\;\simeq\;
Y^+(\widehat{\mathfrak{gl}}_1).
\]
The full affine Yangian and the $\mathcal W_{1+\infty}$ vertex-algebra
action are obtained only after the double, centre, and evaluation
operations in Theorem~\ref{thm:positive-half-boundary}.
\end{theorem}

\begin{proof}
The \v{C}ech nerve of the one-chart cover has one object in every
simplicial degree and all coface maps are identities. Its totalisation
is the chart CoHA. Schiffmann-Vasserot, \emph{Publ. IHES} 118 (2013),
Theorem~9.1, identify that CoHA with the positive half
$Y^+(\widehat{\mathfrak{gl}}_1)$. The remaining assertion is
Theorem~\ref{thm:positive-half-boundary}.
\end{proof}

\end{document}
